\section{What width compression does and does not say about circuits}
\label{sec:product-normal-form}

\begin{corollary}[A finite-description normal form]
\label{cor:product-finite-description}
Suppose a Boolean function on $\{-1,1\}^d$ has an arbitrary finite-width,
depth-$D\ge2$ tanh representation with global signed margin at least one
and product of Frobenius norms at most $P\ge1$. There is a representation
of the same function, with margin at least one, using at most $D$
weighted layers, polynomial width and polynomially bounded individual
weights in $d,D,P$, and product norm at most $8P$.
Here a purely affine function of the augmented input is allowed as the
case of one weighted layer. The weights may moreover be chosen dyadic
with polynomial bit length, at the expense of replacing the product
bound by $16P$. Consequently the function has a Boolean circuit of size
polynomial in $d,D,P$, with no dependence on the original widths,
weights, or their precision.
\end{corollary}

\begin{proof}
Apply Theorem~\ref{thm:pf-width-compression} on the augmented cube with
error $1/4$, and double its readout. This gives width
$C=\operatorname{poly}(d,D,P)$, product at most $P_1=4P$, and signed
margin at least $3/2$. Apply the small-prefix linearization from the
proof of Theorem~\ref{thm:product-target-log-padding}, with
\[
 \tau^2=\frac1{16P_1D(d+1)^{3/2}}.
\]
It changes every output by at most $1/48$. Collapse the resulting linear
blocks and double the final readout. The resulting signed margin exceeds
two, the product is at most $2P_1=8P$, and every physical weight is at
most
\[
 R=\max\{1,2P_1\sqrt C/\tau\}
   =\operatorname{poly}(d,D,P).
\]
The number of weighted layers does not increase. If every hidden
activation is removed, the remaining affine map has norm at most
$2P_1$ and satisfies the same claims.

For completeness, polynomial precision can be imposed without assuming
that the displayed real weights were efficiently computable. Let $D'$
be the remaining depth and choose a polynomial $B\ge2$ that bounds the
Frobenius norm and the row absolute sum of every remaining matrix, and
also bounds all input and hidden dimensions. Each matrix norm $M_j$
satisfies
\[
 M_j\ge\frac1{\sqrt{d+1}\,B^{D'-1}}.
\]
Indeed the margin gives
$1\le |z(x)|\le\sqrt{d+1}\prod_iM_i$, and all other norms are at most
$B$. Round each entry to a dyadic mesh $2^{-b}$, where
\[
 b=O\bigl(D'\log B+\log(d+D'+2)\bigr)
\]
with a sufficiently large universal constant. The matrix Frobenius
perturbations are then at most $M_j/(4D')$, so the product increases
by at most $(1+1/(4D'))^{D'}<2$. The standard layerwise error recursion,
using the row-sum bound $B$ and the one-Lipschitz property of tanh,
makes the total output change less than $1/2$ at the same precision.
The initial margin exceeding two therefore leaves margin at least one.
There are polynomially many entries and $b$ is polynomial in $d,D,P$.

Finally, evaluate each affine map and tanh to sufficiently many binary
digits and propagate a total error smaller than $1/2$. Arithmetic on
these polynomial-length numbers and approximation of tanh on the
polynomially bounded preactivation range use polynomial-size Boolean
circuits. The output sign remains exact because of the margin.
\end{proof}

The converse direction is not supplied by width compression: a
polynomial-size Boolean circuit is not thereby given a polynomial
product-norm tanh representation. In particular, widening an existing
implementation is not itself a proof of such a bound. If a polynomial
product-norm implementation exists at any finite width, the corollary
shows that a polynomial-width, finite-description implementation exists
as well.

Nor does this normal form impose constant threshold depth. The
rank-one approximation argument gives retained depth of order
$P^4\log(2P)$, after restriction to the at most $P^2$ relevant
coordinates, which is only a polynomial depth bound when $P$ is
polynomial. Its combination with width compression therefore gives
no inclusion in $\mathrm{TC}^0$ or $\mathrm{NC}^1$ and no separation of
those circuit classes.
